\chapter{The Cubical-Duck-Type Metatheory}
\label{ch:cubical}

This chapter formalises the metatheory of \Deppy{} as it exists in
the source tree.  The core data lives in:
\begin{itemize}[leftmargin=1.7em]
  \item \texttt{deppy/core/types.py} --- types and terms;
  \item \texttt{deppy/core/kernel.py} --- judgments, proof terms,
    \texttt{ProofKernel.verify};
  \item \texttt{deppy/core/cubical\_effects.py} ---
    $\Diamond$/$\Box$ modalities and effect fibrations;
  \item \texttt{deppy/core/path\_engine.py} ---
    \texttt{PathConstructor}, \texttt{TransportEngine},
    \texttt{UnivalenceEngine};
  \item \texttt{deppy/core/path\_compression.py} --- the
    homotopy-equivalent-branch reduction.
\end{itemize}

We give the syntax, the typing/verification rules, and the trust
discipline.  We are explicit where the metatheory \emph{admits}
without proof.

\section{Syntax: types and terms}
\label{sec:cubical-syntax}

\subsection{Types}

The universe of \Deppy{} types is generated by the \texttt{SynType}
hierarchy in \texttt{deppy/core/types.py}.  We summarise the
constructors we use throughout the rest of the book:

\begin{align*}
  T \;{::=}\;& \mathsf{PyObj} \mid \mathsf{int} \mid \mathsf{float} \mid \mathsf{bool} \mid \mathsf{str} \mid \mathsf{None}\\
            &{}\mid \mathsf{list}[T] \mid \mathsf{dict}[T_k, T_v] \mid \mathsf{tuple}[T_1, \dots, T_n] \mid \mathsf{set}[T] \\
            &{}\mid \mathsf{class}\langle n; \overline{(m : T)}\rangle \\
            &{}\mid \Pi(x : T_1).\, T_2(x) \mid \Sigma(x : T_1).\, T_2(x) \\
            &{}\mid t_1 =_T t_2     \quad \text{(path / identity)}\\
            &{}\mid \mathsf{U}_n     \quad \text{(universe at level $n$)}\\
            &{}\mid \{x : T \mid \varphi(x)\}     \quad \text{(refinement)}\\
            &{}\mid T_1 \cup T_2     \quad \text{(union)}\\
            &{}\mid \mathbb{I}     \quad \text{(cubical interval)}\\
            &{}\mid \mathsf{Glue}(\dots) \\
            &{}\mid \Box_e T \mid \Diamond_e T     \quad \text{(effect-modal)}\\
            &{}\mid \mathsf{Protocol}\langle \overline{(m : T)}\rangle \\
            &{}\mid \mathsf{Optional}[T] \mid \mathsf{Interval}[a, b]
\end{align*}

The Python primitives (\texttt{int}, \texttt{float}, \texttt{bool},
\dots, \texttt{class}, \texttt{tuple}, \texttt{list}) are
\emph{nominal} carriers; the HoTT formers ($\Pi$, $\Sigma$, identity,
universe, refinement) are \emph{structural}.  The cubical layer
($\mathbb{I}$, \texttt{Glue}, $\Box$/$\Diamond$) is structural
\emph{and} explicit about effects.

\subsection{Terms}

The term grammar from \texttt{deppy/core/types.py}:

\begin{align*}
  t \;{::=}\;& x  \quad \text{(variable)}\\
            &{}\mid \overline{c} \quad \text{(literal)}\\
            &{}\mid \lambda x{:}T.\, t \mid t_1\,t_2\\
            &{}\mid \langle t_1, t_2\rangle \mid \mathsf{fst}\,t \mid \mathsf{snd}\,t\\
            &{}\mid \langle\lambda i.\, t\rangle \mid t\,@\,r \quad \text{(path abs/app)}\\
            &{}\mid \mathsf{transport}_T(p, t) \mid \mathsf{comp}_T(\dots) \\
            &{}\mid \mathsf{glue}(\dots) \mid \mathsf{unglue}(\dots) \\
            &{}\mid \mathsf{let}\,x = t_1\,\mathsf{in}\,t_2\\
            &{}\mid \mathsf{ite}(t_b, t_t, t_e)\\
            &{}\mid \mathsf{pyCall}(t_f, \overline{t_a}) \mid \mathsf{getAttr}(t, m) \mid \mathsf{getItem}(t, t_k)
\end{align*}

The four \texttt{pyCall}/\texttt{getAttr}/\texttt{getItem}/\texttt{ite}
constructors are deliberately surface-Python --- they let the kernel
talk about Python-shaped expressions without forcing a translation
into HoTT primitives.  This is the \emph{operational} concession
that allows \Deppy{} to work on idiomatic Python code.

\subsection{Contexts and judgments}

The \texttt{Context} is a list of typed bindings $x_1 : T_1, \dots,
x_n : T_n$ (file: \texttt{deppy/core/types.py}).  A
\texttt{Judgment} is one of:
\begin{align*}
  \Gamma \vdash T : \mathsf{Type}     &\quad\text{(\texttt{WELL\_FORMED})}\\
  \Gamma \vdash t : T                  &\quad\text{(\texttt{TYPE\_CHECK})}\\
  \Gamma \vdash t_1 \equiv t_2 : T     &\quad\text{(\texttt{EQUALITY})}\\
  \Gamma \vdash \mathsf{satisfies}(t, T) &\quad\text{(refinement check)}\\
  \Gamma \vdash \mathsf{covers}(\overline{T_i}, T) &\quad\text{(fibration coverage)}
\end{align*}

The kernel's \texttt{verify} method takes a \texttt{Context}, a
\texttt{Judgment}, and a \texttt{ProofTerm}, and returns a
\texttt{VerificationResult} carrying:
\begin{itemize}[leftmargin=1.7em]
  \item \texttt{success : bool};
  \item \texttt{trust\_level : TrustLevel};
  \item \texttt{message : str} (diagnostic);
  \item \texttt{axioms\_used : list[str]} (provenance).
\end{itemize}

\section{Trust levels}
\label{sec:trust-levels}

The kernel ranks accepted proofs on a seven-step scale (file:
\texttt{deppy/core/kernel.py}, lines 40--65):

\begin{center}
\begin{tabular}{rl}
\toprule
6 & \texttt{KERNEL}            \\
5 & \texttt{Z3\_VERIFIED}      \\
4 & \texttt{STRUCTURAL\_CHAIN} \\
3 & \texttt{LLM\_CHECKED}      \\
2 & \texttt{AXIOM\_TRUSTED}    \\
1 & \texttt{EFFECT\_ASSUMED}   \\
0 & \texttt{UNTRUSTED}         \\
\bottomrule
\end{tabular}
\end{center}

The semantics:
\begin{description}
  \item[\texttt{KERNEL}] every sub-proof was accepted by an internal
    rule.  This is the strongest level reachable without external
    discharge.
  \item[\texttt{Z3\_VERIFIED}] one or more sub-proofs were
    discharged by an SMT call.  We trust the SMT solver's answer.
  \item[\texttt{STRUCTURAL\_CHAIN}] the proof is a structural
    composition of accepted steps but contains some \texttt{Structural}
    or \texttt{Assume} leaves whose semantics is left to the user.
  \item[\texttt{LLM\_CHECKED}] the kernel could verify the structure
    of the proof but at least one obligation was discharged via
    computational testing (i.e.\ running the function on samples).
  \item[\texttt{AXIOM\_TRUSTED}] one or more sidecar axioms were
    cited.  We trust the user's axiomatic claim.
  \item[\texttt{EFFECT\_ASSUMED}] effect modality witnesses were
    accepted without proof (the \texttt{EffectWitness} term).
  \item[\texttt{UNTRUSTED}] no useful verification was performed.
\end{description}

The \texttt{min\_trust} helper (lines 63--65) computes the minimum
trust over a list, giving the standard ``trust composes by minimum''
rule.

\section{Proof terms: the inventory}
\label{sec:proofterm-inventory}

The kernel admits the following \texttt{ProofTerm} subclasses (file:
\texttt{deppy/core/kernel.py}; we group them logically, with the
``$\to$'' showing the kernel rule that fires).  This is the
authoritative reference for \PSDL's compilation targets.

\subsection{Equality and substitution}

\begin{description}
  \item[$\Refl(t)$] proves $t =_T t$ for any $t : T$.
  \item[$\Sym(p)$] from $p : a =_T b$ derive $\Sym(p) : b =_T a$.
  \item[$\Trans(p, q)$] from $p : a =_T b$ and $q : b =_T c$ derive
    $a =_T c$.
  \item[$\Cong(f, p)$] from $p : a =_T b$ derive $f(a) =_{T'} f(b)$.
  \item[$\Transport(F, p, d)$] from $p : a =_T b$ and $d : F(a)$
    derive $F(b)$.
  \item[$\Unfold(\mathit{name}, p)$] reduce the named definition and
    continue with $p$.
  \item[$\Rewrite(L, p)$] rewrite using a propositional equality $L$
    and continue with $p$.
  \item[$\Cut(h, T_h, p_h, p_b)$] introduce hypothesis $h : T_h$
    proved by $p_h$, then prove the body with $p_b$.
  \item[$\mathrm{Assume}(h)$] use a previously bound hypothesis $h$.
\end{description}

\subsection{Induction}

\begin{description}
  \item[$\NatI(x, p_0, p_s)$] $\mathsf{Nat}$-induction on $x$:
    $p_0$ is the zero case, $p_s$ the successor case.
  \item[$\ListI(x, p_n, p_c)$] list induction on $x$.
  \item[{$\Cases(s, [(\mathit{pat}_i, p_i)])$}] case analysis on the
    scrutinee $s$.
\end{description}

\subsection{Cubical / paths}

\begin{description}
  \item[$\PathComp(p, q)$] composition of paths via Kan composition.
  \item[$\Ap(f, p)$] congruence on paths --- $\Ap(f, p) :
    f(a) =_T f(b)$ if $p : a = b$.
  \item[$\FunExt(x, p)$] from a pointwise proof
    $\forall x.\, f(x) = g(x)$ derive $f = g$.
  \item[{$\KanFill(\mathit{dim}, [\mathit{faces}], \dots)$}] supply the
    missing face of a cubical box.
  \item[{$\Glue([\mathit{patches}], [\mathit{overlaps}])$}] glue local
    proofs on a cover into a global proof; requires explicit overlap
    cocycle agreement.
  \item[$\Univ(e, A, B)$] from $e : A \simeq B$ derive a path
    $A =_{\mathsf{Type}} B$ (univalence).
  \item[$\Cocycle(\mathit{level}, \overline{c}, \mathit{boundary},
    \dots)$] a $k$-cocycle in the safety cochain complex.
\end{description}

\subsection{Duck typing / dispatch}

\begin{description}
  \item[$\DuckPath(A, B, \overline{(m, p_m)})$] proves $A \simeq B$
    over a protocol's method set, with one method-wise proof $p_m$
    per method $m$.
  \item[$\CDuckPath(\mathit{cond}, e, A, B)$] a duck path that holds
    only when a runtime predicate \emph{cond} does.
  \item[$\FibLook(\mathit{attr}, p_b, \mathit{thru})$] models
    \texttt{\_\_getattr\_\_}-based proxies: a $\Sigma$-projection on
    the attribute \emph{attr} via the delegate \emph{thru}.
  \item[$\Patch(C, m, v', p_c)$] a monkey-patch path
    $C =_{\mathsf{PyClass}} C[m := v']$, valid only when the
    behavioural-contract proof $p_c$ holds.
\end{description}

\subsection{Effect / exception fibres}

\begin{description}
  \item[$\Effect(e, p)$] a witness that effect $e$ is admitted by
    proof $p$.  Always at trust \texttt{EFFECT\_ASSUMED}.
  \item[$\CEffect(\mathit{cond}, e, p, \mathit{tgt})$] a conditional
    effect witness: under precondition \emph{cond}, effect $e$ is
    absent at target \emph{tgt}.  Trust composes from $p$.
\end{description}

\subsection{Structural and SMT}

\begin{description}
  \item[$\Struct(\mathit{description})$] \emph{admits} the obligation
    on the user's authority.  Always at \texttt{STRUCTURAL\_CHAIN}
    trust.  The soundness gate (\Cref{ch:soundness}) refuses
    certification when an verify-block proof has a \texttt{Structural}
    leaf in a non-vacuous position.
  \item[$\Axiomic(\mathit{name}, \mathit{module})$] cite a registered
    axiom or foundation.  Trust \texttt{AXIOM\_TRUSTED}.
  \item[$\Zthr(\mathit{formula}, \mathit{binders})$] discharge the
    formula via the kernel's \Zthree{} backend.  Trust
    \texttt{Z3\_VERIFIED} when SMT returns \texttt{unsat} on the
    negation.
  \item[$\LeanP(\mathit{theorem}, \mathit{proof}, \dots)$] discharge
    by invoking \Lean{}.  Trust \texttt{KERNEL} when \Lean{}
    accepts.
\end{description}

\subsection{Safety / termination terms}

The remaining \texttt{ProofTerm} subclasses (see
\texttt{deppy/core/kernel.py} from line $\sim$580 onward) carry
the safety pipeline's witnesses: \texttt{SourceDischarge},
\texttt{SemanticSafetyWitness}, \texttt{ModuleSafetyComposition},
\texttt{TerminationObligation}, \texttt{SafetyObligation}.  These
are used by the safety-pipeline pass (\texttt{deppy.pipeline.safety\_pipeline});
they are orthogonal to the \PSDL{} verify-block flow but share the
trust-level discipline.

\section{Verification rules}
\label{sec:rules}

We give the formal rules for the eight non-trivial proof-term
constructors.  Notation: $\Gamma \vdash p : J / L$ means ``proof term
$p$ verifies the judgment $J$ at trust level $L$''.

\paragraph{Refl, Sym, Trans, Cong.}
\begin{mathpar}
  \inferrule{\Gamma \vdash a : T}{\Gamma \vdash \Refl(a) : a =_T a / \mathtt{KERNEL}}
  \and
  \inferrule{\Gamma \vdash p : a =_T b / L}{\Gamma \vdash \Sym(p) : b =_T a / L}
\end{mathpar}
\begin{mathpar}
  \inferrule{
    \Gamma \vdash p : a =_T b / L_p \\
    \Gamma \vdash q : b =_T c / L_q
  }{
    \Gamma \vdash \Trans(p, q) : a =_T c / \min(L_p, L_q)
  }
  \and
  \inferrule{
    \Gamma \vdash f : T \to T' \\
    \Gamma \vdash p : a =_T b / L
  }{
    \Gamma \vdash \Cong(f, p) : f(a) =_{T'} f(b) / L
  }
\end{mathpar}

\paragraph{Cut.}
\begin{mathpar}
  \inferrule{
    \Gamma \vdash p_h : T_h / L_h \\
    \Gamma, h : T_h \vdash p_b : J / L_b
  }{
    \Gamma \vdash \Cut(h, T_h, p_h, p_b) : J / \min(L_h, L_b)
  }
\end{mathpar}

\paragraph{Fiber.}
\begin{mathpar}
  \inferrule{
    \Gamma \vdash s : \mathsf{PyObj} \\
    \forall i.\ \Gamma, s : T_i \vdash p_i : J / L_i \\
    \Gamma \vdash \mathsf{covers}(\overline{T_i}, \mathsf{type}(s))
  }{
    \Gamma \vdash \Fiber(s, [(T_i, p_i)]) : J / \min_i L_i
  }
\end{mathpar}

The \texttt{covers} judgment is decided by checking that the
\emph{exhaustive} flag is set and the union of $T_i$ contains
$\mathsf{type}(s)$ in the type registry; otherwise the kernel
adds an implicit ``else: \texttt{Structural}('uncovered')'' branch
that downgrades trust to \texttt{STRUCTURAL\_CHAIN}.

\paragraph{Z3.}
\begin{mathpar}
  \inferrule{
    \mathtt{z3.solve}(\neg\varphi[\overline{x : T}]) = \mathtt{unsat}
  }{
    \Gamma \vdash \Zthr(\varphi, \overline{(x, T)})
      : \{\mathsf{True} \mid \varphi\} / \mathtt{Z3\_VERIFIED}
  }
\end{mathpar}

When the SMT call returns \texttt{sat} the kernel \emph{rejects}
the proof and additionally surfaces the model as a counterexample.
\texttt{unknown} is \emph{not} accepted; the kernel returns failure.

\paragraph{LeanProof.}
\begin{mathpar}
  \inferrule{
    \mathtt{lean}(\texttt{theorem T : } \varphi \texttt{ := by } P) \;\rightsquigarrow\; \mathtt{exit\,0}
  }{
    \Gamma \vdash \LeanP(\varphi, P, \dots) : \{\mathsf{True} \mid \varphi\} / \mathtt{KERNEL}
  }
\end{mathpar}

\paragraph{Structural and AxiomInvocation.}
\begin{mathpar}
  \inferrule{}{
    \Gamma \vdash \Struct(d) : J / \mathtt{STRUCTURAL\_CHAIN}
  }
  \and
  \inferrule{
    \mathit{name} \in \mathit{axioms} \\
    \mathit{axioms}[\mathit{name}].\mathit{statement} \approx J
  }{
    \Gamma \vdash \Axiomic(\mathit{name}, m) : J / \mathtt{AXIOM\_TRUSTED}
  }
\end{mathpar}

The ``$\approx$'' approximate-match is essential: the kernel parses
the axiom statement as Python AST and pattern-matches it against the
goal.  Exact textual match is the success case; canonicalising
rewrites recover from minor variable-renaming and operator-form
differences.

\paragraph{Cocycle.}
\begin{mathpar}
  \inferrule{
    \forall i.\ \Gamma \vdash c_i : J_i / L_i \\
    \forall (i, j) \in \mathit{boundary\_pairs}.\ \Gamma \vdash o_{ij} : c_i \approx_{(i \cap j)} c_j / L_{ij} \\
    \Gamma \vdash z : \delta\overline{c} = 0 / L_z\ \text{(when supplied)}
  }{
    \Gamma \vdash \Cocycle(k, \overline{c}, \mathit{bp}, \overline{o}, z) : J / \min_{(i, j, z)}(L_i, L_{ij}, L_z)
  }
\end{mathpar}

The boundary-coherence check is non-trivial: each overlap proof
$o_{ij}$ witnesses that components $c_i$ and $c_j$ \emph{agree on
their overlap}.  When boundary pairs are not supplied the kernel
falls back to a coarser ``each component independently''
verification and downgrades the cocycle's classification to
\emph{unstructured}.

\section{Cubical effects}
\label{sec:cubical-effects}

We sketch the modal layer; the source is
\texttt{deppy/core/cubical\_effects.py}.

\subsection{Effect categories}

The \texttt{EffectCategory} enum is:
\begin{equation*}
  \mathsf{PURE},\, \mathsf{IO},\, \mathsf{STATE},\, \mathsf{EXCEPTION},\,
  \mathsf{ALLOCATION},\, \mathsf{NONDETERMINISM},\, \mathsf{DIVERGENCE}.
\end{equation*}
An \texttt{EffectSignature} carries a set of categories, a set of
``resources'' (file handles, network endpoints), and a set of
``modifies'' (mutated variables).  The \texttt{compose} operation is
union; the \texttt{subsumes} order is component-wise inclusion.

\subsection{Modal types}

For an effect signature $E$:
\begin{align*}
  \Box_E A   &\quad\text{guaranteed-pure $A$ when $E \subseteq \{\mathsf{PURE}\}$}\\
  \Diamond_E A &\quad\text{maybe-effectful $A$ producing values when no exception}
\end{align*}
The expected unit-and-bind structure holds:
\begin{align*}
  \mathsf{lift}     &: A \to \Diamond_E A\\
  \mathsf{bind}     &: \Diamond_E A \to (A \to \Diamond_E B) \to \Diamond_E B\\
  \mathsf{purify}   &: \Box_\mathsf{PURE} A \to A
\end{align*}
The \texttt{CubicalEffectVerifier} (file:
\texttt{deppy/core/cubical\_effects.py}, line 240+) constructs a
real \texttt{TransportProof} when transporting an effect proof
along an equivalence path between two function ASTs.

\subsection{Exception fibres}

An \texttt{EXCEPTION}-effect call to a function $f : A \to B$ is
modelled as having codomain $\Diamond_{\{\mathsf{EXCEPTION}\}} B$,
whose \emph{value fibre} is $B$ and whose \emph{exception fibre}
indexes a Python exception type.  A property
$\varphi : B \to \mathsf{Prop}$ holds of $f$ when:
\begin{itemize}[leftmargin=1.7em]
  \item the value fibre of $f$ has type $\{x : B \mid \varphi(x)\}$, and
  \item the exception fibre is non-empty (contains a witness raise).
\end{itemize}
The kernel records the second condition with \texttt{EffectWitness}
or \texttt{ConditionalEffectWitness}.  The pipeline emits
\texttt{vacuous "..."} (kernel \texttt{Structural}) for the
exception fibre by default; \texttt{certify\_or\_die} accepts the
vacuous claim when no counterexample is found.

\section{Path engine}
\label{sec:path-engine}

\texttt{deppy/core/path\_engine.py} contains the constructors and
verifiers for the path layer:

\begin{itemize}[leftmargin=1.7em]
  \item \texttt{PathConstructor.function\_path(fn1, fn2)} produces a
    \texttt{PathAbs} witness of $fn_1 \simeq fn_2$ when their bodies
    differ only in renaming/algebraic-identity-applications.
  \item \texttt{TransportEngine.transport(motive, p, d)} verifies the
    transport rule and returns a checked \texttt{TransportProof}
    term.
  \item \texttt{UnivalenceEngine.from\_equiv(e, A, B)} constructs a
    \texttt{Univalence} proof term over an equivalence triple.
  \item \texttt{CechDecomposer.decompose(cover)} returns a Čech-style
    decomposition of a property over a class hierarchy --- the basis
    for the \texttt{CechGlue} term.
\end{itemize}

The path engine is wired into the \PSDL{} compiler primarily through
the \texttt{transport(...)} and \texttt{equiv\_to\_path(...)} sugar.
See \Cref{ch:psdl}.

\section{The cubical-duck-type model: pulling it together}
\label{sec:ddtm-pulling-together}

The combined picture is best read off the example.  Take
\texttt{Point.distance(p, q) >= 0}.  We want a proof that for
\emph{every} $p, q$ inhabiting the \emph{actual} \Sympy{}
\texttt{Point} class, the value-fibre of the call satisfies $\ge 0$.

\begin{enumerate}[leftmargin=1.7em]
  \item Resolve \verb|Point.distance| via \Importlib; read its source
    via \Inspect.
  \item Notice the body has shape
    \begin{verbatim}
    if isinstance(other, Point): ...
    else: ... raise TypeError ...
    \end{verbatim}
  \item This is a \emph{fibration} on \texttt{isinstance}: the source
    type is $\mathsf{PyObj}$, the fibres are $\mathsf{Point}$ and
    its complement.
  \item In the \texttt{Point}-fibre, the body returns
    \texttt{sqrt(sum\_zip\_sub\_sq(self, other))}, with
    \texttt{sum\_zip\_sub\_sq} a sum of squares (always non-negative)
    and \texttt{sqrt} non-negative on non-negatives (foundation
    \texttt{Real\_sqrt\_nonneg}).
  \item In the complement fibre, the body raises a
    \texttt{TypeError}: this is the \emph{exception fibre} of the
    $\Diamond$-modality.  The property is vacuously preserved
    because there is no return value to compare.
\end{enumerate}

The corresponding \texttt{ProofTerm} encodes step (3) as a
\texttt{Fiber}, step (4) inside the value-fibre branch as
\texttt{Cut(\dots, \texttt{Z3Proof}(...), \dots, AxiomInvocation("Real\_sqrt\_nonneg"))},
and step (5) in the complement-fibre branch as a
\texttt{Structural("vacuous: exception fibre")}.

The trust level composes as
\[
  \min(\mathtt{Z3\_VERIFIED}, \mathtt{AXIOM\_TRUSTED}, \mathtt{STRUCTURAL\_CHAIN})
  \;=\;\mathtt{STRUCTURAL\_CHAIN}.
\]
The certificate documents this and \texttt{certify\_or\_die}
inspects whether the \texttt{Structural} leaf is in a vacuous
position; in this case it is, so the gate passes.

\section{Soundness, in the only sense we can claim it}
\label{sec:soundness-honest}

We have not given a denotational semantics for \Deppy{} terms, and
we have not proved cut-elimination.  What we \emph{do} claim:

\begin{theorem}[Engineering soundness]
\label{thm:engsound}
For every verify-block whose certificate prints \texttt{CERTIFIED},
the following hold:
\begin{enumerate}[leftmargin=1.7em]
  \item The \texttt{ProofTerm} the kernel built passes the rules of
    \S\ref{sec:rules}.
  \item Every \texttt{Z3Proof} sub-term was checked by \Zthree{} (or
    by the cached lookup of a previous identical check).
  \item Every \texttt{LeanProof} sub-term was checked by \Lean{} (or
    by cache).
  \item Every cited foundation has either been \Zthree-verified or
    been admitted with the foundation declaration visible in the
    certificate.
  \item The certificate's \Lean{} text was round-tripped via
    \texttt{check\_lean\_module\_source} and \Lean{} returned
    \texttt{exit 0}.
  \item The \texttt{certify\_or\_die} list of failure conditions is
    empty.
\end{enumerate}
\end{theorem}

\begin{proof}[Justification]
Items (1)--(5) follow by inspection of the implementation.  Item (6)
is enforced operationally: \texttt{certify\_or\_die} (file
\texttt{deppy/lean/sidecar\_runner.py}, lines 62--98) refuses
\texttt{CERTIFIED} unless every item is satisfied.
\end{proof}

\begin{remark}
Theorem~\ref{thm:engsound} is \emph{not} a logical-consistency
result.  It is a statement about what the implementation does.  A
mathematical-soundness theorem would require modelling Python's
semantics, the SMT solver's semantics, and \Lean's metatheory; we
have done none of those things.

The implication for the user is that the trust surface bottoms at:
\Lean's kernel, \Zthree{}'s answer to a satisfiability query,
CPython's faithful execution of the source bytes that
\Inspect{} returned, and the user's hand-written admissions
(foundations, opaque axioms, structural leaves).  We hope this is a
trust surface a careful user can reason about.
\end{remark}

\section{Summary}
\label{sec:ch4-summary}

We have given the metatheory:
\begin{itemize}[leftmargin=1.7em]
  \item \textbf{Types} include the Python carriers, the HoTT formers,
    and the cubical/effect layer.
  \item \textbf{Terms} include the standard $\lambda$-calculus,
    HoTT path constructors, and surface-Python idioms.
  \item \textbf{Judgments} are the standard four plus
    fibration-coverage.
  \item \textbf{Trust levels} run from \texttt{KERNEL} (full) to
    \texttt{UNTRUSTED} (none); \texttt{min\_trust} composes.
  \item \textbf{Proof terms} are 38+ subclasses spanning equality,
    induction, cubical, duck-type, effect, structural, and external
    discharge primitives.
  \item \textbf{The verification rules} are local on each constructor
    and composable via \texttt{min\_trust}.
  \item \textbf{The cubical effects} make the value/exception
    fibres of Python calls explicit.
  \item \textbf{Soundness} is engineering soundness (Theorem
    \ref{thm:engsound}), not metamathematical soundness.
\end{itemize}

The next chapter introduces the syntax of the sidecar files in which
all this is invoked: the \texttt{.deppy} format.
